\chapter{Model gate: GNN experimental opcional}
\label{sec:paper-gate}

\section{El dilema del componente ML}

Un clasificador ML dentro de un sistema productivo tiene dos modos
de falla sistémicos más allá de su accuracy:

\begin{enumerate}
    \item \textbf{Bloqueo por modelo ausente}. El cliente recién
    instalado no tiene modelo entrenado; si el sistema exige uno,
    no puede ni empezar.
    \item \textbf{Degradación silenciosa}. El modelo produce
    predicciones sin confianza adecuada; el sistema las integra
    como si fueran ciertas; el analista las trata como
    groundtruth.
\end{enumerate}

El patrón \emph{model gate} es la respuesta arquitectónica a ambas.
La funcionalidad ML es opcional por contrato, no por
configuración.

\section{Contrato}

El clasificador vive en \code{core/gnn/} detrás de un trait:

\begin{lstlisting}[language=Rust]
pub trait ThreatClassifier: Send + Sync {
    fn classify_path(&self, path: &[NodeRef], edges: &[EdgeRef]) -> f64;
    fn model_info(&self) -> ModelInfo;
}

pub struct NoModel;
impl ThreatClassifier for NoModel {
    fn classify_path(&self, _: &[NodeRef], _: &[EdgeRef]) -> f64 { 0.0 }
    fn model_info(&self) -> ModelInfo { ModelInfo::none() }
}

pub struct OnnxClassifier { session: OrtSession, ... }
impl ThreatClassifier for OnnxClassifier { ... }
\end{lstlisting}

El composer de scoring no conoce cuál está activo. Recibe un
\code{Arc<dyn ThreatClassifier>}. Si la configuración no apunta a
un modelo ONNX, se instancia \code{NoModel} y devuelve 0.
\code{w\_5} es 0 por defecto: multiplicado por 0 no penaliza.

\section{Arquitectura del modelo v2}

El modelo actual (\emph{v2 experimental}) es una GraphSAGE de dos
capas sobre el path instanciado. Entrada: features por nodo
(type one-hot, grado log-normalizado, first-seen-age) y por
arista (type one-hot, timestamp delta al anterior). Salida:
probabilidad de clase ``threat''.

Entrenamiento fue sobre un dataset etiquetado de logs Sentinel AAD
con incidentes reales etiquetados por un analista senior
($\sim$2000 paths positivos, $\sim$20000 negativos sampleados).
Export a ONNX vía \code{tch-rs}. Carga en runtime con \code{ort}.

\section{Separación layer-wise}

El modelo ML vive en el \emph{core} (\code{core/gnn}), consume
tipos del grafo pero no conoce la superficie HTTP/MCP. El switch
entre \code{NoModel} y \code{OnnxClassifier} está en un módulo de
configuración (\code{core/graph-engine/src/scoring/gate.rs}). La
decisión se documenta en ADR-003.

\section{Honestidad epistémica}

\begin{itemize}
    \item El modelo \emph{no está entrenado} para FortiAnalyzer,
    Cognito ni Zoho. Tirarlo sobre esos datasets sin retraining
    da scores no informativos. Por eso el peso por defecto es 0.
    \item El modelo \emph{no es confiable como único juez}. Sirve
    como un quinto voto en la composición, no como clasificador
    autónomo.
    \item Un modelo por cliente / por vertical es probablemente el
    camino. Esto es una observación, no una decisión tomada.
\end{itemize}

\section{Extensiones previstas}

\begin{itemize}
    \item Exposición del vector de confianza (no sólo el score), útil
    en UI para distinguir ``GNN está seguro'' de ``GNN no
    sabe''.
    \item Training pipeline reproducible (\code{training/})
    ya scaffoldea pero está fuera del alcance de este reporte.
    \item Model gate per-vertical: selección del ONNX por
    cliente/contexto, no global.
\end{itemize}
